\section{Results}\label{sec:results}

We organise the evaluation around four questions. Section~\ref{sec:setup} first
establishes Beacon's accuracy on the held-out 2025 test set
(Table~\ref{tab:headline}, Figure~\ref{fig:horizon}); we then turn to the live,
horizon-matched comparison against ISO-NE's operational forecast
(Table~\ref{tab:iso}); isolate the contributions of the target-hour features and
the unified weather source (Table~\ref{tab:ablation}, Figure~\ref{fig:ablation});
and finally audit the pipeline for data leakage (Table~\ref{tab:leak}). Throughout,
test-set numbers refer to the $7{,}823$ held-out windows from 2025, whereas the
live numbers refer to the $\sim$30-day, $720$-point operational window.

\subsection{Held-out test performance}
Table~\ref{tab:headline} summarises Beacon's accuracy on the 2025 held-out set.
At the easiest horizon ($h=1$) the model attains an MAE of $60.9$\,MW, and at the
day-ahead horizon ($h=24$) an MAE of $76.7$\,MW, for an average of $73.6$\,MW
across all $24$ horizons. The $h=1$ forecast reaches a MAPE of $2.25\%$ and an
$R^2$ of $0.977$, explaining nearly all of the load variance. Crucially, error
grows only mildly with horizon: the day-ahead MAE is just $1.26\times$ the
one-hour MAE, evidence that the direct per-horizon design largely contains the
error compounding that plagues recursive forecasters. Averaged across the $24$
horizons, Beacon improves on the single-model rolled-out baseline by $80.6\%$.

\begin{table}[t]
\centering
\caption{Held-out test-set (2025) headline metrics for Beacon.}
\label{tab:headline}
\begin{tabular}{lc}
\toprule
Metric & Value \\
\midrule
MAE @ $h=1$ & $60.9$ MW \\
MAE @ $h=24$ (day-ahead) & $76.7$ MW \\
Average MAE across $24$\,h & $73.6$ MW \\
MAPE ($h=1$) & $2.25\%$ \\
$R^2$ ($h=1$) & $0.977$ \\
Horizon degradation (MAE$_{h24}$/MAE$_{h1}$) & $1.26\times$ \\
Improvement vs.\ single-model (avg.) & $+80.6\%$ \\
\bottomrule
\end{tabular}
\end{table}

\subsection{Per-horizon behaviour and the single-model collapse}
Figure~\ref{fig:horizon} contrasts the per-horizon MAE of Beacon against the
single-model rolled-out baseline. Beacon stays flat across the entire horizon,
rising only from $60.9$\,MW at $h=1$ to $76.7$\,MW at $h=24$. The single-model
baseline, by contrast, degrades catastrophically at the mid-horizons: from
$62.6$\,MW at $h=1$ it climbs to $234.8$\,MW at $h=4$, $351.7$\,MW at $h=6$, and
peaks near $505.0$\,MW around $h=12$, before partially recovering to $234.8$\,MW
at $h=24$. The mechanism is intuitive: a model trained to predict $t+1$ uses the
load at $t$ as a near-perfect anchor, but reusing that same prediction $12$ hours
later places it at the opposite phase of the diurnal cycle, where the previous
hour is a poor proxy. The recovery near $h=24$ reflects the return of the daily
cycle to a similar phase. Beacon sidesteps this entirely by training a dedicated
model per horizon, each conditioned on the target-hour calendar and weather.

\begin{figure}[t]
\centering
\includegraphics[width=\columnwidth]{fig_horizon_mae.png}
\caption{Per-horizon test-set MAE: Beacon (direct multi-horizon) versus the
single-model rolled-out baseline. Beacon remains flat ($60.9$--$76.7$\,MW),
while the single model collapses at mid-horizons, peaking near $505$\,MW around
$h=12$ where the one-step anchor is least informative.}
\label{fig:horizon}
\end{figure}

\subsection{Live comparison against ISO-NE}
Table~\ref{tab:iso} reports the horizon-matched live evaluation over the
$\sim$30-day, $720$-point window, with both forecasters at the same $24$-hour-ahead
horizon, the same Open-Meteo weather, and ISO-NE real-time demand as ground truth.
Beacon attains an MAE of $89.3$\,MW and a MAPE of $3.50\%$, against ISO-NE's
$92.6$\,MW and $3.70\%$. Beacon thus beats the operational forecast on both MAE
($-3.3$\,MW, $\approx +3.5\%$) and MAPE, while ISO-NE retains a marginally higher
$R^2$ ($0.923$ vs.\ $0.904$). The headline takeaway is that a fully reproducible,
keyless pipeline edges out a mature operational benchmark on absolute error once
the comparison is made honestly horizon-matched. We emphasise that the live window
is short and falls in spring/early-summer, so this result is a like-for-like
demonstration rather than a claim of dominance across all seasons.

\begin{table}[t]
\centering
\caption{Live, horizon-matched day-ahead comparison ($\sim$30-day window,
$720$ points). Both forecasters at $24$\,h ahead, same Open-Meteo weather,
ISO-NE real-time demand as actuals.}
\label{tab:iso}
\begin{tabular}{lccc}
\toprule
Forecaster & MAE (MW) & MAPE (\%) & $R^2$ \\
\midrule
Beacon (day-ahead) & $\mathbf{89.3}$ & $\mathbf{3.50}$ & $0.904$ \\
ISO-NE (day-ahead)  & $92.6$ & $3.70$ & $\mathbf{0.923}$ \\
\bottomrule
\end{tabular}
\end{table}

\subsection{Ablations}
Table~\ref{tab:ablation} decomposes the day-ahead gains into two cumulative
changes, with Figure~\ref{fig:ablation} visualising the day-ahead ($h=24$) MAE
across the three configurations. Starting from a lags-only model served
OpenWeatherMap (OWM) weather, the day-ahead MAE is $183.0$\,MW (with $h=1$ at
$78.7$\,MW). Adding the target-hour exogenous features---the forecasted weather
and exactly-known calendar at hour $t+h$---cuts the day-ahead MAE to $110.3$\,MW
(a $\approx 40\%$ reduction) while leaving $h=1$ essentially unchanged at
$77.7$\,MW, as expected since the nowcast already has its target-hour context.
Finally, unifying the weather source on Open-Meteo for both training and serving
drops the day-ahead MAE to $76.7$\,MW and, notably, improves $h=1$ to $60.9$\,MW.

\begin{table}[t]
\centering
\caption{Cumulative ablation of target-hour features and weather source
(test-set MAE, MW).}
\label{tab:ablation}
\begin{tabular}{lcc}
\toprule
Configuration & MAE @ $h=1$ & MAE @ $h=24$ \\
\midrule
Lags only, OWM weather & $78.7$ & $183.0$ \\
\;+ target-hour features, OWM & $77.7$ & $110.3$ \\
\;+ Open-Meteo (matched) --- \textbf{final} & $\mathbf{60.9}$ & $\mathbf{76.7}$ \\
\bottomrule
\end{tabular}
\end{table}

\begin{figure}[t]
\centering
\includegraphics[width=\columnwidth]{fig_ablation.png}
\caption{Day-ahead ($h=24$) test-set MAE across the three cumulative
configurations: lags-only ($183.0$\,MW) $\rightarrow$ $+$ target-hour weather
($110.3$\,MW) $\rightarrow$ $+$ Open-Meteo train/serve matching ($76.7$\,MW).}
\label{fig:ablation}
\end{figure}

The deployment consequence of the weather source is sharper still. A model trained
on OpenWeatherMap CSVs but served Open-Meteo weather suffers a live day-ahead MAE
of $\approx 147$\,MW; retraining on Open-Meteo to match the serving source brings
this down to $\approx 90$\,MW---essentially the ISO-NE level ($\approx 92$\,MW).
Matching the train/serve weather source alone thus removes most of the gap to the
operational benchmark, underscoring that a silent train/serve distribution
mismatch can nearly double deployed error.

\subsection{Leakage audit}
Finally, Table~\ref{tab:leak} reports the shuffled-target test, the gold standard
for detecting feature--target leakage. With genuine targets the test $R^2$ is
$0.956$ at $h=1$ and $0.770$ at $h=24$; when the targets are randomly permuted
before training, $R^2$ collapses to $-0.019$ and $-0.021$ respectively---no better
than predicting the mean. This confirms the model exploits real feature--target
structure rather than any inadvertent leakage. The monotone increase of MAE with
horizon documented in Figure~\ref{fig:horizon} is corroborating evidence: a leaking
model would remain near-flat across horizons rather than degrade in the orderly,
physically plausible way Beacon does.

\begin{table}[t]
\centering
\caption{Leakage audit: real-target versus shuffled-target test $R^2$.}
\label{tab:leak}
\begin{tabular}{lcc}
\toprule
Horizon & Real-target $R^2$ & Shuffled-target $R^2$ \\
\midrule
$h=1$  & $0.956$ & $-0.019$ \\
$h=24$ & $0.770$ & $-0.021$ \\
\bottomrule
\end{tabular}
\end{table}
